% Generated draft: requires mathematical and source review.
\documentclass[11pt]{article}
\usepackage[margin=25mm]{geometry}
\usepackage{fontspec}
\setmainfont{Latin Modern Roman}
\usepackage{amsmath,amssymb,mathtools,unicode-math}
\setmathfont{Latin Modern Math}
\usepackage{enumitem,needspace,xurl,hyperref}
\setlength{\parindent}{0pt}
\setlength{\parskip}{5pt}
\setlength{\emergencystretch}{3em}
\newcommand{\N}{\mathbb{N}}
\newcommand{\Z}{\mathbb{Z}}
\newcommand{\Q}{\mathbb{Q}}
\newcommand{\R}{\mathbb{R}}
\newcommand{\C}{\mathbb{C}}
\newcommand{\F}{\mathbb{F}}
\newcommand{\A}{\mathbb{A}}
\newcommand{\PP}{\mathbb{P}}
\newcommand{\E}{\mathbb{E}}
\newcommand{\eps}{\varepsilon}
\newcommand{\dd}{\,\mathrm{d}}
\begin{document}
\section*{UnsolvedMath ID: 30003163}
\textbf{Source problem code:} catalogue entry
\subsection{Definitions and mathematical statement}
Let \(K\) be a complete discretely valued field, let \(\mathcal O_K\) be its valuation ring, and let \(\mathfrak m_K\) be the maximal ideal of \(\mathcal O_K\). Assume that the residue field \(k=\mathcal O_K/\mathfrak m_K\) is finite and has characteristic \(p>0\). Let \(D\) be a central division algebra over \(K\), so that its center is \(K\). The source does not explicitly state finite-dimensionality over \(K\), although that is the usual convention for central division algebras in this reduced-norm setting.

Write \(\mathbf K(A)\) for the algebraic \(K\)-theory spectrum of a ring \(A\), and \(K_j(A)=\pi_j\mathbf K(A)\) for its integral \(K\)-groups in the degrees under consideration. Retain the source notation \(K_j(A,\mathbb Z_p)\) for its \(p\)-adic \(K\)-groups. The source does not specify the degree range, whether connective or nonconnective \(K\)-theory is intended, or the precise convention for these coefficients.

\textbf{Source question, with unresolved qualifications.} For every pair \((K,D)\) satisfying these hypotheses, the record asks whether the maps it calls reduced-norm isomorphisms
\[
\operatorname{Nrd}_{D/K}\colon K_j(D,\mathbb Z_p)\xrightarrow{\;\sim\;}K_j(K,\mathbb Z_p)
\]
are induced by a map of \(K\)-theory spectra. An integral lifting interpretation would ask for a map
\[
f\colon\mathbf K(D)\longrightarrow\mathbf K(K)
\]
whose induced maps on \(p\)-adic \(K\)-groups agree with the displayed reduced norms in the intended degrees. The source also asks for an integral reduced-norm isomorphism on \(K\)-groups. This would require corresponding isomorphisms
\[
\pi_j(f)\colon K_j(D)\xrightarrow{\;\sim\;}K_j(K)
\]
if the integral isomorphisms are intended to arise from the same map. However, the source does not distinguish a map between completed spectra from a map between integral spectra, and it does not justify its asserted equivalence between spectrum-level realization and integral isomorphisms. These requirements are therefore recorded as related questions, without asserting their equivalence.

For an abelian group \(G\), define its \(p\)-primary torsion subgroup by
\[
G[p^\infty]=\{x\in G:\text{there exists }n\geq 0\text{ with }p^n x=0\}.
\]
This subgroup has bounded exponent precisely when
\[
\exists e\geq 0\quad\forall x\in G[p^\infty],\qquad p^e x=0.
\]
The record additionally asks whether the \(p\)-primary torsion subgroups of the integral \(K\)-groups have bounded exponent. It does not specify which of \(K_j(K)\) and \(K_j(D)\) are intended, the relevant degrees, or whether an exponent may depend on the field, algebra, and degree. Moreover, its supplied literature assessment already describes the needed torsion as known to have bounded exponent. This conflict prevents identifying a single precise unresolved torsion assertion from the supplied evidence.
\subsection{Short English statement}
For a central division algebra over a complete discretely valued field with finite residue field, the source asks whether its stated \(p\)-adic reduced-norm isomorphisms arise from a map of algebraic \(K\)-theory spectra and extend to integral isomorphisms. The spectrum interpretation, degree range, and accompanying torsion question require clarification; the claimed equivalence of the formulations is not established by the supplied record.
\subsection{Sources}
\begin{itemize}
\item Local UnsolvedMath registry: \texttt{data/unsolvedmath/problems.json}, permanent ID \(30003163\), supplied complete source record, title \emph{Integral Reduced Norms in Division-Algebra K-Theory}, exact source problem code \texttt{catalogue entry}. No individual proposer is supplied.
\item Supplied report citation: \emph{Algebraic K-theory and Motivic Cohomology (2017)}, Oberwolfach Reports, pp.~1753--1807, DOI \texttt{10.4171/owr/2016/31}. \url{https://doi.org/10.4171/owr/2016/31}. The record describes Oberwolfach Reports as documenting mathematical workshops and their open-problem discussions; it separately records proposed year \(2016\).
\item Supplied literature reference and attribution: Hesselholt--Larsen--Lindenstrauss, \emph{p-adic K-theory}. \url{https://arxiv.org/abs/1807.00104}. The supplied assessment attributes a proof of the \(p\)-adic analogue to Hesselholt, Larsen, and Lindenstrauss.
\item Supplied literature reference: \emph{Oberwolfach Report 31}. \url{https://doi.org/10.4171/owr/2016/31}.
\item Literature assessment dated exactly \texttt{2026-08-21}: the supplied curation labels the problem open, says the integral reduced-norm isomorphism remains unknown, and explains this by saying that the needed \(p\)-primary torsion is only known to have bounded exponent. This is supplied dated curation derived from the OWR record, not independently verified literature assessment. The source itself requests independent verification of time-sensitive status and cited claims.
\item The registry labels the statement \texttt{corrected\_verified} and reports a manual semantic, typographic, and LaTeX review. Those are supplied registry qualifications, not independent verification of the mathematical formulation or the asserted equivalence.
\end{itemize}
\end{document}
